\begin  {titlepage}
\null
\begin{flushright}
 {\it  Version: November 2000 } 
\end{flushright}
\vfill
%
\title  { \huge\bf  LatMRG \medskip \\ \LARGE\bf User's Guide\\ \bigskip
          \large\bf A Modula-2 software for the theoretical analysis \\
            of linear congruential and multiple recursive \\
            random number generators }

\author { Pierre L'Ecuyer \ and \ Raymond Couture \\ \medskip
          \small\rm  % GERAD \\ and \\
          D\'epartement d'Informatique et de Recherche Op\'erationnelle \\
          Universit\'e de Montr\'eal, C.P. 6128, Succ.\ Centre-Ville,
          Montr\'eal, Canada, H3C 3J7 }

{\noindent\large\bf Abstract. \ }
{\tt LatMRG} is a software toolkit for examining theoretical properties of
linear congruential or multiple recursive random number generators.
It is implemented as a library of modules written in Modula-2.
It offers tools to check whether a generator has maximal period or not,
to apply the lattice and spectral tests (in large dimensions),
and to perform computer searchs for good (or bad) generators 
according to different criteria.
One can analyse the lattice structure of points formed by successive
values in the generator's sequence, or formed by ``leapfrog'' values.
Generators with large moduli and multipliers (e.g. numbers of many
hundreds of bits), as well as combined generators, can also be analyzed.
Multiply-with-carry generators can also be studied by analyzing their
corresponding linear congruential generators. 

\keywords {{\large\bf Keywords:} 
            Random number generation; linear congruential generators;
            multiple recursive generators;  multiply-with-carry;
            lattice structure; spectral test; maximal period }

\end {titlepage}
